% ===========================================================================
\section{Two explicit constructions: Shannon and Fano}\label{sec:shannon-fano}
\warmth{0}\quad\whisper{Kraft only promised that a code exists; here are two recipes that build one.}

\begin{strand}
Proposition~\ref{prop:shannon-code} produced the lengths
$\lceil-\log_Dp(x)\rceil$ but not the codewords. Shannon's recipe reads them off
the cumulative distribution function; Fano's recipe splits the alphabet into
balanced halves. Shannon's version also reveals the price of coding with the
\emph{wrong} distribution: exactly a relative entropy.
\end{strand}

\trigger{Use this section when you need an actual codeword, or when the code was designed for a distribution other than the true one.}

\block{Shannon's code}

The idea is to give $x$ a codeword of length
$\ell_x=\lceil\log_D\frac1{p(x)}\rceil$, and to choose \emph{which} string of
that length by cutting the $D$-ary expansion of the cumulative probability.

\begin{definition}[Shannon code]\label{def:shannon-code-construction}
Let $\X=\{x_1,\ldots,x_m\}$ with pmf $p$ and let $|\D|=D$.
\begin{enumerate}
  \item Order the probabilities in \answerin[3cm]{decreasing} order,
  $p_1\ge p_2\ge\cdots\ge p_m$.
  \item Set the cumulative probability
  $\alpha_k=\sum_{i=1}^{k-1}p_i$, and define $C(x_k)$ to be the first
  \[
    \ell_k=\answerin[2.6cm]{\lceil\log_D(1/p_k)\rceil}
  \]
  digits in the $D$-ary expansion of $\alpha_k$.
\end{enumerate}
\studentrecall{Sort decreasingly; cut the $D$-ary expansion of $\alpha_k=\sum_{i<k}p_i$ after $\lceil\log_D(1/p_k)\rceil$ digits.}
\end{definition}

The truncation error satisfies $0\le\alpha_k-0.C(x_k)<D^{-\ell_k}\le p_k$, and
$\alpha_{k+1}-\alpha_k=p_k$, so distinct codewords describe disjoint
subintervals of $[0,1)$. Hence the code is a prefix code with exactly the
lengths used in Proposition~\ref{prop:shannon-code}.

\begin{yourturn}
Why does the construction need the probabilities sorted, and where would the
argument fail without it?
\studentrecall{Sorting keeps the intervals in order so a short codeword never sits inside a later one.}
\ifshowanswers\else\workspace[2]\fi
\answerprose{Sorting makes $\ell_1\le\ell_2\le\cdots$, so the interval attached
to a later symbol is shorter than the truncation resolution of every earlier
one; no codeword can then be a prefix of a previous codeword. Without sorting a
long, low-probability interval could contain the truncation point of a later
symbol and prefix-freeness could fail.}
\end{yourturn}

\block{Coding with the wrong distribution}

Suppose we build the Shannon code from a pmf $q$ while the source really follows
$p$. The excess cost is exactly the relative entropy $\KL_D(p\,\|\,q)$.

\begin{proposition}[mismatched Shannon code]\label{prop:mismatched-shannon}
Let $p$ and $q$ be pmfs on $\X$, let $X\sim p$, and let $\D$ be a finite code
alphabet with $|\D|=D$. If $C$ is the Shannon code built from the distribution
$q$, then
\[
  \Ent_D(X)+\KL_D(p\,\|\,q)
  \le\E\bigl[|C(X)|\bigr]
  <\Ent_D(X)+\KL_D(p\,\|\,q)+1 .
\]
\end{proposition}

\begin{proof}
By construction $|C(x)|=\lceil\log_D\frac1{q(x)}\rceil$, so
\[
\begin{aligned}
  \E\bigl[|C(X)|\bigr]
  &=\sum_xp(x)\left\lceil\log_D\frac1{q(x)}\right\rceil\\
  &<\sum_xp(x)\left(\log_D\frac1{q(x)}+1\right)\\
  &=\sum_xp(x)\left(\log_D\frac{p(x)}{q(x)}
    +\log_D\frac1{p(x)}\right)+1\\
  &=\KL_D(p\,\|\,q)+\Ent_D(X)+1 .
\end{aligned}
\]
The lower bound follows from the same chain with the ceiling bounded below by
its argument, i.e.\ by deleting the $+1$.
\end{proof}

\begin{yourturn}
In the chain above, which single algebraic step creates the divergence term?
\studentrecall{Insert $p(x)/p(x)$ inside the logarithm to split $\log_D\frac1{q}$ into $\log_D\frac{p}{q}+\log_D\frac1p$.}
\ifshowanswers\else\workspace[2]\fi
\answerprose{Writing $\frac1{q(x)}=\frac{p(x)}{q(x)}\cdot\frac1{p(x)}$ and
splitting the logarithm. The first piece averages to $\KL_D(p\|q)$, the second
to $\Ent_D(X)$.}
\end{yourturn}

\begin{yourturn}
Recover Proposition~\ref{prop:shannon-code} from
Proposition~\ref{prop:mismatched-shannon}.
\studentrecall{Take $q=p$, so $\KL_D(p\|p)=0$.}
\ifshowanswers\else\workspace[1]\fi
\answerprose{Take $q=p$: then $\KL_D(p\|p)=0$ and the bounds collapse to
$\Ent_D(X)\le\E[|C(X)|]<\Ent_D(X)+1$.}
\end{yourturn}

\block{Fano's code (1949)}

Fano's construction is top-down: repeatedly split the alphabet into two groups
of as nearly equal probability as possible, and let the split decide the next
code digit.

\begin{definition}[Fano code, binary case]\label{def:fano-code}
Let $p$ be a pmf on $\X=\{x_1,\ldots,x_m\}$ and $X\sim p$.
\begin{enumerate}
  \item Sort the symbols so that $p_1\ge p_2\ge\cdots\ge p_m$.
  \item Find the index $r$ minimizing
  \[
    \left|\sum_{i\le r}p_i-\sum_{i>r}p_i\right|,
  \]
  and split $\X$ into the two groups
  $\X_0=\{x_i:i\le r\}$ and $\X_1=\{x_i:i>r\}$.
  \item Define the first digit of every codeword in $\X_0$ to be
  \answerin[1cm]{$0$}, and in $\X_1$ to be \answerin[1cm]{$1$}.
  \item Repeat steps (2) and (3) inside each group, appending one digit per
  round, until every group is a singleton.
\end{enumerate}
\studentrecall{Split into two halves of near-equal probability; the side you land on is the next digit.}
This produces a code $C_F:\X\to\{0,1\}^*$.
\end{definition}

Because each symbol corresponds to a leaf of the splitting tree, $C_F$ is a
\answerin[2.4cm]{prefix} code, and one can show
\[
  \E\bigl[|C_F(X)|\bigr]\le\Ent_D(X)+\answerin[0.9cm]{2}.
\]
\studentrecall{Prefix code; expected length at most $\Ent_D(X)+2$.}
The bound is one symbol weaker than Shannon's, and indeed \emph{Fano's code is
not optimal in general}: balancing probabilities greedily at each split can be
beaten by Huffman's bottom-up construction.

\begin{yourturn}
Rank the three codes of this lecture --- Shannon, Fano, optimal (Huffman) --- by
the guarantee each carries on $\E[|C(X)|]$, and say which of them is actually
optimal.
\studentrecall{$\Ent_D<\Ent_D+1$ (Shannon) $\le\Ent_D+2$ (Fano); only Huffman is optimal.}
\ifshowanswers\else\workspace[3]\fi
\answerprose{All three are bounded below by $\Ent_D(X)$. Shannon's code
guarantees $<\Ent_D(X)+1$, Fano's only $\le\Ent_D(X)+2$, and the optimal code
is at least as good as both, so it also satisfies $<\Ent_D(X)+1$. Neither
Shannon's nor Fano's code is optimal in general; the optimum is attained by the
Huffman code.}
\end{yourturn}

\vspace{0.8cm}
\recall{Which quantity measures the loss from designing a code for $q$ when the source is $p$?}
